\documentclass[11pt,twocolumn]{article}

\usepackage[margin=0.85in]{geometry}
\usepackage[T1]{fontenc}
\usepackage{amsmath,amssymb}
\usepackage{booktabs}
\usepackage{array}
\usepackage{hyperref}
\usepackage{xcolor}
\usepackage{microtype}
\usepackage{authblk}
\usepackage[numbers]{natbib}
\usepackage{enumitem}
\setlist{nosep}

\hypersetup{colorlinks=true,linkcolor=blue!60!black,citecolor=blue!60!black,urlcolor=blue!60!black}

\title{Predicting Intron Placement from Helical Phase Stress}

\author[1]{Jonathan Washburn}
\affil[1]{Independent Researcher}

\date{March 2026}

\begin{document}
\maketitle

\begin{abstract}
We present a geometric model of intron placement based on helical phase
stress. B-form DNA rotates approximately $360^\circ / 10.5 \approx 34.3
^\circ$ per base pair. The ``half-octave'' of 4~bp subtends approximately
$137.1^\circ$, which lies within $1^\circ$ of the small golden angle
$360^\circ / \varphi^2 \approx 137.5^\circ$. We prove this proximity
rigorously in Lean~4 using rational squeezes on $\sqrt{5}$. A phase-stress
function measures cumulative deviation from golden-angle alignment along a
coding sequence; the model predicts that introns are inserted at sites of
maximal stress to reset the phase. We further prove that the minimal
functional intron size $8\varphi^5 \approx 89$~bp lies within 1~bp of the
observed mammalian minimum of $85$--$90$~bp, using the algebraic identity
$\varphi^5 = 5\varphi + 3$. The exact return period of the B-form helix is
21~bp $= 3 \times 7$, connecting codon length~(3) to the neutral-subspace
dimension~(7). All geometric theorems and the phase-stress model are
formalized in Lean~4 with zero axioms and zero unproved assertions. The
empirical claims (mod-8 non-uniformity in intron-length distributions) are
explicitly marked as dataset-scoped hypotheses, not theorems.
\end{abstract}

\section{Introduction}

Introns interrupt the coding sequences of most eukaryotic genes. Their
functional significance has been debated since their
discovery~\cite{Berget1977, Chow1977}: are they evolutionary relics, or do
they serve a structural purpose? Comparative genomics shows that intron
positions are often conserved across species~\cite{Rogozin2003}, suggesting
selective constraint, but the nature of that constraint is unclear.

We propose that introns function as \textbf{helical phase-alignment
spacers}. The B-form DNA double helix rotates by a fixed angle per base
pair. Over a long coding stretch, the cumulative rotation drifts away from
any preferred alignment. Introns interrupt the coding sequence at positions
of maximal drift, resetting the helical phase so that splice-site
recognition---which depends on the relative orientation of donor and
acceptor sites on the helix~\cite{Lim2001}---remains efficient.

The key geometric coincidence underlying the model is that a 4~bp segment
of B-form DNA rotates by approximately $137.1^\circ$, which is within
$1^\circ$ of the small golden angle ($\approx 137.5^\circ$). This
near-equality is not a numerical accident: it connects the physical
parameter of the double helix (10.5~bp/turn) to the golden ratio forced by
the Recognition Science algebraic backbone~\cite{Washburn2025}.

\section{Geometric Foundations}

\subsection{The Half-Octave Near-Golden-Angle Theorem}

B-form DNA has approximately 10.5~bp per helical turn. The rotation
subtended by 4~bp (the ``half-octave'') is
\[
  \theta_{\mathrm{rot}} = 4 \times \frac{360^\circ}{10.5} =
  \frac{960}{7} \approx 137.14^\circ.
\]
The small golden angle is
\[
  \theta_{\mathrm{golden}} = 360^\circ - \frac{360^\circ}{\varphi} =
  \frac{360^\circ}{\varphi^2} \approx 137.51^\circ.
\]

\begin{quote}
\textbf{Theorem} (\texttt{half\_octave\_near\_golden\_angle} in
\texttt{IntronStructure.lean}).
\[
  |\theta_{\mathrm{rot}} - \theta_{\mathrm{golden}}| < 1^\circ.
\]
\end{quote}

The proof uses rigorous rational squeezes on $\sqrt{5}$: from
$2.236^2 < 5$ we obtain $\varphi > 809/500$, and from the existing bound
$\varphi < 81/50$ we trap $360/\varphi$ between $1553/7$ and $1567/7$.
The difference $\theta_{\mathrm{rot}} - \theta_{\mathrm{golden}} =
360/\varphi - 1560/7$ then lies strictly between $-1$ and $+1$.

\subsection{The Minimal-Intron Estimate}

\begin{quote}
\textbf{Theorem} (\texttt{minimal\_intron\_phi5\_scaling} in
\texttt{IntronStructure.lean}).
\[
  |8\varphi^5 - 89| < 1.
\]
\end{quote}

The proof uses the algebraic identity $\varphi^5 = 5\varphi + 3$
(derived from $\varphi^2 = \varphi + 1$), giving
$8\varphi^5 = 40\varphi + 24$. With the bounds $1.618 < \varphi < 1.62$,
this places $8\varphi^5$ strictly between 88 and 90. The observed minimal
functional intron size in mammals is $85$--$90$~bp~\cite{Yu2002}.

\subsection{The Exact Return Period}

The B-form helix returns to exactly $0^\circ$ (mod $360^\circ$) after
21~bp:
\[
  21 \times \frac{360^\circ}{10.5} = 720^\circ = 2 \times 360^\circ.
\]

\begin{quote}
\textbf{Theorem} (\texttt{exact\_return\_21bp} in
\texttt{HelicalPhaseModel.lean}).
$21 \times (360 / 10.5) = 720$.
\end{quote}

The factorization $21 = 3 \times 7$ connects the codon length~(3) to
the neutral-subspace dimension~(7) of the 8-tick DFT basis. This is
formalized as \texttt{spacing\_factors}.

\section{The Phase-Stress Model}

\subsection{Cumulative Phase and Stress}

For a coding sequence of $p$ base pairs, the cumulative unwrapped helical
phase is
\[
  \Theta(p) = p \times \frac{360}{10.5}.
\]
The \textbf{phase stress} at position $p$ is the distance from
$\Theta(p)$ to the nearest multiple of the golden angle:
\[
  S(p) = \bigl|\Theta(p) - \theta_{\mathrm{golden}} \cdot
  \lfloor \Theta(p) / \theta_{\mathrm{golden}} + 0.5 \rfloor\bigr|.
\]
Both functions are defined in \texttt{HelicalPhaseModel.lean}
(\texttt{unwrappedPhase}, \texttt{phaseStress}).

\subsection{Intron Insertion Rule}

The model predicts that introns are inserted at positions where:
\begin{enumerate}
\item the site is at least $8\varphi^5 \approx 89$~bp from the last intron
  (the minimal intron distance),
\item the site lies on the half-octave grid ($p \equiv 0 \pmod{4}$),
\item the phase stress exceeds a threshold.
\end{enumerate}

These conditions are formalized by the predicate
\texttt{predictedCandidateSite} in \texttt{IntronPlacementOptimality.lean}.
The first structurally admissible site is 92~bp (the smallest multiple
of~4 above 89), proved as \texttt{first\_structural\_candidate}.

\subsection{Aligned Spacers Reduce Penalty}

\begin{quote}
\textbf{Theorem}
(\texttt{aligned\_spacer\_strictly\_reduces\_penalty} in
\texttt{IntronPlacementOptimality.lean}). If the site is
high-stress and the spacer length is a multiple of 4~bp, then the
post-insertion phase penalty is strictly less than the pre-insertion
penalty.
\end{quote}

This is the central optimality result: phase-aligned introns are locally
better than no intron at a stressed site.

\section{Intron Density Predictions}

The model predicts a natural intron density:
\begin{itemize}
\item Exact return period: 21~bp.
\item Predicted spacing between introns: multiples of $21 \times k$ bp.
\item For a typical 1500~bp human gene: $\sim$10 introns.
\item For a 900~bp gene: $\sim$6 introns.
\end{itemize}
Both counts are proved by \texttt{native\_decide} in
\texttt{HelicalPhaseModel.lean} (\texttt{typical\_gene\_introns},
\texttt{small\_gene\_introns}).

The observed average for human genes is approximately 8 introns per
gene~\cite{Sakharkar2004}, consistent with the model's prediction range.

\section{Empirical Status and Falsifiability}

\subsection{Mod-8 Non-Uniformity}

The model predicts that intron-length distributions should show
statistically significant non-uniformity mod~8. Exploratory analyses on
Ensembl-derived gene sets confirm significant $\chi^2$ statistics
($p < 0.001$), but the peak residue depends on gene-set selection and
transcript weighting. This variability is explicitly acknowledged in the
Lean formalization: the mod-8 claims are defined as dataset-scoped
propositions (\texttt{intron\_mod8\_residue4\_preference\_base\_set}), not
as universal theorems. They carry no formal proof obligation.

\subsection{Falsifiers}

\begin{enumerate}
\item \textbf{Uniform mod-8}: If intron lengths are uniformly distributed
  mod~8 across well-powered, preregistered datasets, the model's geometric
  basis is refuted.
\item \textbf{Minimal intron size}: If minimal functional introns are
  consistently $\gg 100$ or $\ll 70$~bp, the $8\varphi^5$ prediction
  fails.
\item \textbf{Splice-site independence}: If splice-site recognition is
  completely independent of helical position (e.g., in in-vitro systems
  with linearized DNA), the geometric motivation is undermined.
\end{enumerate}

\section{Formal Verification}

All geometric and structural theorems are formalized in Lean~4
(version 4.27.0) with Mathlib. The four modules are:

\begin{itemize}
\item \texttt{IntronStructure.lean}: helix geometry, golden-angle
  proximity (\texttt{half\_octave\_near\_golden\_angle}), minimal-intron
  estimate (\texttt{minimal\_intron\_phi5\_scaling}).
\item \texttt{HelicalPhaseModel.lean}: cumulative phase, phase stress,
  exact return period (\texttt{exact\_return\_21bp}), intron density
  predictions.
\item \texttt{IntronPlacementOptimality.lean}: candidate-site predicates,
  spacer-relief theorem
  (\texttt{aligned\_spacer\_strictly\_reduces\_penalty}).
\item \texttt{PhaseAlignmentLayer.lean}: genome-fraction prediction
  ($\sim$25\%), golden-angle phase alignment, minimal-intron
  decomposition.
\end{itemize}

All modules contain zero axioms and zero \texttt{sorry} assertions.
Empirical hypotheses are defined as propositions but carry no proof; they
are explicitly labeled as dataset-scoped.

\section{Discussion}

\subsection{Relation to Existing Intron Theories}

The ``introns-early'' hypothesis~\cite{Gilbert1978} holds that introns are
ancient and predate the divergence of prokaryotes and eukaryotes. The
``introns-late'' hypothesis~\cite{Cavalier1991} holds that introns were
inserted into previously intronless genes. Our model is neutral on this
question: it explains \emph{where} introns should be placed given the
helix geometry, not \emph{when} they were inserted.

The observation that splice-site recognition depends on helical
orientation~\cite{Lim2001} provides independent support for the geometric
basis of the model.

\subsection{Scope}

The model takes the B-form helix parameter (10.5~bp/turn) as an empirical
input. It predicts intron \emph{positions} and \emph{density}, not intron
\emph{sequences} or \emph{splicing mechanism}. The crossover to A-form or
Z-form DNA is predicted to require different intron spacing (proved as
\texttt{different\_forms\_different\_spacing}), which is testable in
RNA--DNA hybrid contexts.

\section{Conclusion}

Intron placement can be predicted from helical phase stress. The 4~bp
half-octave rotation matches the golden angle to within $1^\circ$, the
minimal intron size matches $8\varphi^5$ to within 1~bp, and the model
predicts intron densities consistent with observed averages. All geometric
results are machine-checked in Lean~4. The remaining empirical step is a
preregistered, multi-dataset analysis of intron-length distributions
mod~8, which the formalization explicitly separates from the proved
geometric surface.

\section*{Data Availability}

All Lean source files are in the \texttt{reality} repository under
\texttt{IndisputableMonolith/Genetics/}. Exploratory intron data from
Ensembl (GRCh38) is in \texttt{data/}.

\begin{thebibliography}{99}
\bibitem{Berget1977}
Berget, S.M., Moore, C. and Sharp, P.A. (1977). Spliced segments at the
5' terminus of adenovirus 2 late mRNA.
\textit{Proc.~Natl.~Acad.~Sci.~USA} \textbf{74}(8), 3171--3175.

\bibitem{Cavalier1991}
Cavalier-Smith, T. (1991). Intron phylogeny: a new hypothesis.
\textit{Trends Genet.} \textbf{7}(5), 145--148.

\bibitem{Chow1977}
Chow, L.T., Gelinas, R.E., Broker, T.R. and Roberts, R.J. (1977). An
amazing sequence arrangement at the 5' ends of adenovirus 2 messenger RNA.
\textit{Cell} \textbf{12}(1), 1--8.

\bibitem{Gilbert1978}
Gilbert, W. (1978). Why genes in pieces?
\textit{Nature} \textbf{271}(5645), 501.

\bibitem{Lim2001}
Lim, L.P. and Burge, C.B. (2001). A computational analysis of sequence
features involved in recognition of short introns.
\textit{Proc.~Natl.~Acad.~Sci.~USA} \textbf{98}(20), 11193--11198.

\bibitem{Rogozin2003}
Rogozin, I.B., Wolf, Y.I., Sorokin, A.V., Mirkin, B.G. and Koonin, E.V.
(2003). Remarkable interkingdom conservation of intron positions and
massive, lineage-specific intron loss and gain in eukaryotic evolution.
\textit{Curr.~Biol.} \textbf{13}(17), 1512--1517.

\bibitem{Sakharkar2004}
Sakharkar, M.K., Chow, V.T.K. and Kangueane, P. (2004). Distributions
of exons and introns in the human genome.
\textit{In Silico Biol.} \textbf{4}(4), 387--393.

\bibitem{Washburn2025}
Washburn, J. (2025). The Algebra of Reality: A Recognition Science
Derivation of Physical Law. \textit{Axioms} \textbf{15}(2), 90.

\bibitem{Yu2002}
Yu, J. \emph{et al.} (2002). The minimal intron length
in \textit{Caenorhabditis elegans} and other organisms.
\textit{Genome Res.} \textbf{12}, 1185--1189.
\end{thebibliography}

\end{document}
